\chapter{Built-in Datasets and Data Visualization in R}
\label{chap:visualization}

\chapterguide
  {You should be able to load data frames and inspect their columns.}
  {load built-in datasets; create pie charts, box plots, histograms, scatterplots, bar charts, and line graphs; add labels, titles, legends, colors, and multiple series using the exact families of functions shown in Lab 8a.}

\section{Using built-in datasets}

The handout begins with \texttt{data()} to list available sample datasets and then loads \texttt{mtcars}:

\begin{lstlisting}[style=dsr]
data()

data(mtcars)
print(head(mtcars, 10))

input <- mtcars[, c("mpg", "cyl")]
print(head(input))
\end{lstlisting}

This creates a simple workflow for visualization practice: load a known dataset, inspect a subset, and pass selected columns into plotting functions.

\section{Built-in graphics demonstrations}

\begin{lstlisting}[style=dsr]
demo(graphics)
demo(image)
\end{lstlisting}

The Lab 8a instruction notes that pressing Enter advances through the demonstrations.

\section{Pie charts}

The basic pie chart uses a vector of values and labels:

\begin{lstlisting}[style=dsr]
x <- c(21, 62, 10, 53)
labels <- c("London", "New York", "Singapore", "Mumbai")

pie(x, labels)
\end{lstlisting}

The handout then saves the chart to a PNG file:

\begin{lstlisting}[style=dsr]
png(file = "city.png")
pie(x, labels)
dev.off()
\end{lstlisting}

\subsection{Title, colors, percentages, and legend}

\begin{lstlisting}[style=dsr]
pie(
  x,
  labels,
  main = "City pie chart",
  col = rainbow(length(x))
)

piepercent <- round(100 * x / sum(x), 1)

pie(
  x,
  labels = piepercent,
  main = "City pie chart",
  col = rainbow(length(x))
)

legend(
  "topright",
  c("London", "New York", "Singapore", "Mumbai"),
  cex = 0.8,
  fill = rainbow(length(x))
)
\end{lstlisting}

The percentage expression explicitly converts each value into its share of the total before plotting.

\subsection{3D pie chart}

The handout introduces \texttt{plotrix} and \texttt{pie3D()}:

\begin{lstlisting}[style=dsr]
library(plotrix)

# As printed in the handout:
pie3D(
  x,
  labels = lbl,
  explode = 0.1,
  main = "Pie Chart of Countries"
)
\end{lstlisting}

\begin{pitfall}
The Lab 8a 3D pie-chart fragment uses \texttt{lbl}, but no \texttt{lbl} object is defined in the supplied sequence. A student must define the desired labels (for example from the earlier \texttt{labels} vector) before this fragment can run as written.
\end{pitfall}

\section{Box plots}

The handout uses a formula to compare mileage across cylinder counts:

\begin{lstlisting}[style=dsr]
boxplot(
  mpg ~ cyl,
  data = mtcars,
  xlab = "Number of Cylinders",
  ylab = "Miles Per Gallon",
  main = "Mileage Data"
)
\end{lstlisting}

A more customized version adds notches, variable widths, colors, and supplied group names:

\begin{lstlisting}[style=dsr]
boxplot(
  mpg ~ cyl,
  data = mtcars,
  xlab = "Number of Cylinders",
  ylab = "Miles Per Gallon",
  main = "Mileage Data",
  notch = TRUE,
  varwidth = TRUE,
  col = c("green", "yellow", "purple"),
  names = c("High", "Medium", "Low")
)
\end{lstlisting}

\section{Histograms}

\begin{lstlisting}[style=dsr]
v <- c(9, 13, 21, 8, 36, 22, 12, 41, 31, 33, 19)

hist(
  v,
  xlab = "Weight",
  col = "yellow",
  border = "blue"
)
\end{lstlisting}

The next handout fragment adds axis limits and a requested break count. Its printed argument contains a spacing/wording issue:

\begin{lstlisting}[style=dsr]
# Handout text shows: n breaks = 5
# Working form for the intended argument:
hist(
  v,
  xlab = "Weight",
  col = "green",
  border = "red",
  xlim = c(0, 40),
  ylim = c(0, 5),
  breaks = 5
)
\end{lstlisting}

\begin{pitfall}
The source prints \texttt{n breaks = 5}, which is not a valid named argument in that form. The listing above clearly labels \texttt{breaks = 5} as the working interpretation rather than silently reproducing the typo.
\end{pitfall}

\section{Scatterplots}

Lab 8a selects weight and mileage from \texttt{mtcars}:

\begin{lstlisting}[style=dsr]
input <- mtcars[, c("wt", "mpg")]

plot(
  x = input$wt,
  y = input$mpg,
  xlab = "Weight",
  ylab = "Milage",
  xlim = c(2.5, 5),
  ylim = c(15, 30),
  main = "Weight vs Milage"
)
\end{lstlisting}

A scatterplot matrix expands the idea to four variables:

\begin{lstlisting}[style=dsr]
pairs(
  ~ wt + mpg + disp + cyl,
  data = mtcars,
  main = "Scatterplot Matrix"
)
\end{lstlisting}

\section{Bar charts}

The simplest bar chart is:

\begin{lstlisting}[style=dsr]
H <- c(7, 12, 28, 3, 41)
barplot(H)
\end{lstlisting}

The customized version supplies month names, labels, title, fill color, and border color:

\begin{lstlisting}[style=dsr]
H <- c(7, 12, 28, 3, 41)
M <- c("Mar", "Apr", "May", "Jun", "Jul")

barplot(
  H,
  names.arg = M,
  xlab = "Month",
  ylab = "Revenue",
  col = "blue",
  main = "Revenue chart",
  border = "red"
)
\end{lstlisting}

\subsection{Grouped / stacked-style matrix input}

Lab 8a constructs a 3-by-5 matrix of values and passes it to \texttt{barplot()} with a region legend:

\begin{lstlisting}[style=dsr]
colors <- c("green", "orange", "brown")
months <- c("Mar", "Apr", "May", "Jun", "Jul")
regions <- c("East", "West", "North")

Values <- matrix(
  c(2,9,3,11,9, 4,8,7,3,12, 5,2,8,10,11),
  nrow = 3,
  ncol = 5,
  byrow = TRUE
)

barplot(
  Values,
  main = "total revenue",
  names.arg = months,
  xlab = "month",
  ylab = "revenue",
  col = colors
)

legend("topleft", regions, cex = 1.3, fill = colors)
\end{lstlisting}

\section{Line graphs}

\begin{lstlisting}[style=dsr]
v <- c(7, 12, 28, 3, 41)
plot(v, type = "o")
\end{lstlisting}

The handout then adds title, labels, and color:

\begin{lstlisting}[style=dsr]
plot(
  v,
  type = "o",
  col = "red",
  xlab = "Month",
  ylab = "Rain fall",
  main = "Rain fall chart"
)
\end{lstlisting}

A second series is added using \texttt{lines()}:

\begin{lstlisting}[style=dsr]
t <- c(14, 7, 6, 19, 3)

plot(
  v,
  type = "o",
  col = "red",
  xlab = "Month",
  ylab = "Rain fall",
  main = "Rain fall chart"
)
lines(t, type = "o", col = "blue")
\end{lstlisting}

\begin{dsfigure}
\begin{tikzpicture}[
  box/.style={draw=DSBlue,fill=DSPaleBlue,rounded corners=2mm,text width=3.0cm,minimum height=1.05cm,align=center,font=\small\bfseries\color{DSNavy}},
  arr/.style={-{Stealth[length=2mm]},thick,draw=DSTeal}
]
\node[box] (question) {What do you want to show?};
\node[box,below left=1.2cm and 2.6cm of question] (part) {Parts of a whole\\Pie};
\node[box,below=1.2cm of question] (dist) {Distribution / groups\\Histogram, box plot};
\node[box,below right=1.2cm and 2.6cm of question] (rel) {Relationship / trend\\Scatter, line};
\draw[arr] (question) -- (part);
\draw[arr] (question) -- (dist);
\draw[arr] (question) -- (rel);
\end{tikzpicture}
\end{dsfigure}

The diagram is only a study organizer for the plot families present in the handout; it does not replace the requirement to explain why a particular chart supports the observation you are making.

\begin{quickcheck}
\begin{enumerate}
  \item Which command lists built-in datasets and which command loads \texttt{mtcars}?
  \item How are percentages calculated for the pie chart in the handout?
  \item Which plot compares \texttt{mpg} across \texttt{cyl} groups?
  \item What does \texttt{pairs()} produce in the supplied example?
  \item Which function adds a second series to an existing line plot?
\end{enumerate}
\end{quickcheck}

\begin{chapterrecap}
\begin{itemize}
  \item Lab 8a covers six plot families: pie, box, histogram, scatter, bar, and line.
  \item Plot customization includes titles, axis labels, legends, colors, limits, and multiple series.
  \item Built-in datasets such as \texttt{mtcars} provide ready-made material for practicing visual analysis.
  \item The handout also demonstrates saving a chart to PNG and using package-specific 3D pie charts.
  \item A visualization in Lab 8b is expected to support an observation or story, not merely exist as decoration.
\end{itemize}
\end{chapterrecap}
